% Visual companion through checkpoint 49. All figures are Python-generated.
\begingroup
\renewcommand{\atlasfigure}[2]{\begin{center}\includegraphics[width=.97\linewidth,height=.415\textheight,keepaspectratio]{figures/#1.pdf}\end{center}{\small\color{muted}#2}\par}

\section{How to read this atlas}

The development moves from finite algebra to an infinite-dimensional operator. Each passage changes the question that must be answered. A matrix may preserve the norm of one spinor without defining an operator on every square-summable field. A finite Fourier calculation may be exact without controlling the omitted modes. A symmetric formula still needs a specified domain and an adjoint-domain argument before self-adjointness is established.

The following twelve plates isolate these transitions. They follow a single set of conventions: the unit torus $\T^3=(\R/\Z)^3$, normalized Haar measure, Fourier momentum $p(n)=2\pi n$, and four-component column spinors with coordinates $0,1,2,3$. Time evolution uses $U(t)=\exp(-itH)$ on each finite matrix. The centered symbol is
\[
q_h(n)_j=\frac{\sin(2\pi h n_j)}{h}\quad(h\ne0),
\qquad q_0(n)=2\pi n.
\]
This is a Fourier multiplier on the same torus. The integer modes have not been identified modulo a sampled grid.

\paragraph{Three kinds of visual evidence.}
The projector tree and transport square explain definitions and relationships. Their geometry is schematic. The matrix display records exact entries. Dispersion, evolution, and tail plots are recomputed numerical illustrations of explicit formulas. A smooth curve does not certify a theorem, and a finite partial sum does not decide convergence of an infinite series. The accompanying text supplies the analytic distinction in each case.

\paragraph{The complete free-field argument.}
The operator construction supplies full coefficient and physical evolution, a self-adjoint Hamiltonian, finite Fourier graph approximation, and the exact first-order domain. The checked generator connects the group to this operator. The centered extension proves graph and strong-resolvent limits and the quadratic rate under a finite third-order Fourier moment. The abstract spectral-functional-calculus identification remains a separate written deduction.

\paragraph{Reproduction.}
The retained figure generator is \path{reference_record/atlas_tools/make_figures.py}. The script checks the Clifford relations in exact symbolic arithmetic and compares the displayed propagator with an independent matrix-exponential calculation on a stated finite set. The script specifies its software dependencies, tolerances and finite ranges, and writes a manifest named \path{evidence/figure_checks.json} in its original tree. That manifest is not present in this publication copy; these atlas computations were not rerun for the current formal extension. The current scope schematic is separately regenerated by \path{tools/render_formal_scope.py} from an accepted formal receipt. These checks validate the examples used by the atlas. They do not substitute for the Lean source or its verification receipts.

\clearpage
\section{Refinement preserves an exact decomposition}
\atlasfigure{01_refinement}{A four-dimensional coordinate example. Each child is an orthogonal projector and the children sum to their parent. The drawing is a schematic of refinement, not an identification with spatial dyadic boxes.}

Take $P=\operatorname{diag}(1,1,0,0)$ and $Q=I-P$. The first split distinguishes two orthogonal sectors. Refining $P$ into its first two coordinate projectors preserves its range exactly. Refining $Q$ does the same for the remaining coordinates:
\[
P=P_0+P_1,\qquad Q=Q_2+Q_3,\qquad
I=P_0+P_1+Q_2+Q_3.
\]
An operator $T$ can be resolved into transfer blocks $P_aTP_b$. A block records which input sector is accepted and which output sector is retained. Summing all blocks recovers $T$ when the projectors resolve the identity. Off-diagonal blocks describe transfer between distinct sectors, while diagonal blocks describe action within one sector.

The value of the language is its bookkeeping. It distinguishes an exact decomposition from an approximation that discards some blocks. This elementary example also explains why a change of representation must preserve both the projectors and the inner product when it is used to transport norm claims.


\clearpage
\section{A Fourier mode lives on the three-torus}
\atlasfigure{02_torus_modes}{Three genuine three-dimensional slices of the real part of $e_{(1,1,1)}(x)$. Equal values on opposite faces reflect periodicity. Color shows a real component, not probability density.}

The scalar mode $e_n(x)=\exp(2\pi i n\cdot x)$ has absolute value one everywhere. For normalized Haar measure, its squared $L^2$ norm is also one. Distinct integer modes are orthogonal:
\[
\int_{\T^3}\overline{e_n(x)}e_k(x)\,dx=\delta_{nk}.
\]
The plate shows the mode $(1,1,1)$ at $x_3=0,1/2,1$. Its real part changes sign, while the magnitude of the full complex wave stays fixed. The upper and lower faces have the same pattern because adding one to a coordinate changes the phase by an integer multiple of $2\pi$.

A spinor amplitude $a_n\in\C^4$ multiplies this scalar mode. For a finite sum $f(x)=\sum_n e_n(x)a_n$, orthogonality removes the cross terms only after integration. It gives $\norm{f}_{L^2}^2=\sum_n\norm{a_n}^2$. Local interference can still produce a nonconstant density. This distinction is essential when reading the evolution plate.


\clearpage
\section{The coordinate convention determines the signs}
\atlasfigure{03_clifford}{Real and imaginary parts of the four supplied $4\times4$ generators. Blank cells are zero. The lower panel records coefficients of $i$, so a displayed $-1$ denotes the entry $-i$.}

Rows and columns follow the spinor order $0,1,2,3$. In particular,
\[
\alpha_yv=(-iv_3,\;iv_2,\;-iv_1,\;iv_0)^T.
\]
The sign in this formula fixes the worked packet's fourth component. Replacing it silently with a different convention would change the explicit evolution even if another Clifford representation remained mathematically valid.

The matrices are Hermitian involutions and anticommute in distinct directions. Therefore the modal Hamiltonian
\[
H(m,p)=m\beta+p_x\alpha_x+p_y\alpha_y+p_z\alpha_z
\]
satisfies $H(m,p)^2=(m^2+|p|^2)I$. The cross terms vanish in pairs. This single identity explains the energy shell and the trigonometric formula for the matrix exponential. Hermiticity supplies exact norm preservation.

The figure generator checks all four squares and all six distinct anticommutators in exact Gaussian-integer arithmetic. That small computation is a reproducibility check of the displayed matrices, separate from the formal proofs in the project.


\clearpage
\section{Dispersion changes when momentum becomes a sine}
\atlasfigure{04_dispersion}{Continuum and centered axis energies at $m=1$, $h=1/12$. Curves interpolate real momentum. Dots mark actual integer torus modes on the positive and negative centered branches. The right panel enlarges the low momentum energy difference.}

The square of the Hamiltonian is scalar, so its two shell magnitudes are
\[
E(p)=\sqrt{m^2+|p|^2},\qquad
E_h(p)=\sqrt{m^2+\sum_j[\sin(hp_j)/h]^2}.
\]
Each nonzero shell has positive and negative spectral sectors. The centered symbol is close to $p$ when $|hp_j|$ is small. Over a larger real momentum interval it bends and returns, whereas the continuum shell continues to grow. These global curves should not be read as uniform approximation over all modes.

For the actual torus axis mode $n=(1,0,0)$, $p=(2\pi,0,0)$ and $q_{1/12}=(6,0,0)$. The corresponding positive energies are $\sqrt{1+4\pi^2}$ and $\sqrt{37}$. The torus mode labels remain integers throughout. A curve connecting those labels is a visual interpolation, not an additional family of torus eigenfunctions.


\clearpage
\section{Quadratic consistency has a fixed-mode scope}
\atlasfigure{05_symbol_error}{The centered sine momentum approaches the continuum momentum. At the fixed axis mode $n=(1,0,0)$, the numerical error follows the analytic quadratic bound as $h$ decreases. The sample spacings are $2^{-2},\ldots,2^{-11}$.}

The elementary sine remainder gives
\[
\left|\frac{\sin(hp_j)}{h}-p_j\right|
\leq \frac{|h|^2|p_j|^3}{6}.
\]
For a fixed finite set of modes, the momentum factors have a finite maximum. This yields a quadratic consistency estimate with a constant depending on that set. Exact Hermitian evolution then converts a modal Hamiltonian error into an evolution error growing at most linearly with $|t|$.

The dependence on momentum matters. The factor $|p_j|^3$ grows without bound as the integer mode increases. One cannot remove a Fourier cutoff and retain the same quadratic constant for every $L^2$ field. Additional weighted regularity is needed to sum a rate estimate of this form over infinitely many coefficients.

The checked full-space convergence result instead controls a chosen finite core and its remaining tail separately. The next two plates show why exact stability makes that passage possible without assigning a universal quadratic rate to arbitrary data.


\clearpage
\section{Local density moves while total norm is conserved}
\atlasfigure{06_packet_evolution}{The exact three-mode centered packet with $m=1$, $h=1/12$. Left: pointwise density on the line $x_2=x_3=0$. Right: integrated component weights and their conserved total. A line through the torus is a display slice, not the integration domain.}

Use modes $(0,0,0)$, $(1,0,0)$ and $(1,1,1)$, with initial spinors respectively $e_0/\sqrt2$, $e_0/2$ and $e_0/2$. Their coefficient norm squares sum to one. Each coefficient evolves by
\[
U_n(t)=\cos(E_nt)I-i\frac{\sin(E_nt)}{E_n}H_n,
\]
with the continuous zero-energy branch interpreted as $I$.

The field is a coherent sum of these evolving modes. Cross terms redistribute the local density, and individual spinor components exchange weight. Fourier orthogonality and modal unitarity nevertheless give $\norm{f_h(t)}_{L^2}^2=1$ for every real time. Component 1 remains zero in this particular example and is omitted from the right-hand legend.

The atlas checks the formula against an independent matrix exponential and verifies the total norm on 401 times in $[0,1]$. The all-time statement comes from the algebraic and Fourier arguments, not from the sampled residual.


\clearpage
\section{Removing the cutoff requires two choices}
\atlasfigure{07_two_tail}{An axis-supported example $a_n=n^{-2}/\sqrt{\zeta(4)}$ for $n\geq1$. Left: coefficient mass and a selected cutoff. Right: explicit bounds at $|t|\leq1$ for three fixed cutoffs. The horizontal floors are the two-tail term.}

Write $P_Ka$ for a finite Fourier truncation. The two exact evolutions preserve the norm, so
\[
\norm{(U_h(t)-U(t))a}
\leq\norm{(U_h(t)-U(t))P_Ka}
       +2\norm{(I-P_K)a}.
\]
The correct order is to choose $K$ so that the tail is small, then choose $h$ so that the finite-core term is small uniformly over the time interval. The plot uses the smaller of the finite-core quadratic bound and the elementary stability bound $2\norm{P_Ka}$.

Holding $K$ fixed leaves a positive tail floor even as $h\to0$. Increasing $K$ lowers that floor but increases the coefficient in the finite-core estimate. This tradeoff is not a failure of convergence. It explains why a density-and-stability proof establishes convergence without automatically providing one quadratic rate for all data.

The selected axis sequence is only an illustration. The checked density argument applies to arbitrary square-summable coefficient fields, with finite supports chosen for those fields.


\clearpage
\section{Square summability is weaker than operator membership}
\atlasfigure{08_domain_thresholds}{Finite partial sums through $N=4096$ for axis-supported coefficients $a_{(n,0,0)}=n^{-s}e_0$, $n\geq1$. The zero mode and every other mode are zero. The exact classifications below use the power-series criterion, not the apparent flattening of a curve.}

For this one-dimensional support inside $\Z^3$, ordinary square summability and first-order weighted summability have different thresholds:
\[
\sum_{n\geq1}n^{-2s}<\infty\iff s>\tfrac12,
\qquad
\sum_{n\geq1}(1+4\pi^2n^2)n^{-2s}<\infty
\iff s>\tfrac32.
\]
The choice $s=0.4$ does not define an $L^2$ coefficient field at all. The choice $s=1$ defines an $L^2$ field outside the Dirac domain. The borderline $s=1.5$ still lies outside that domain because its derivative-energy series contains a harmonic tail. The choice $s=2$ belongs to the domain.

These are thresholds for the stated axis-supported family. They must not be transferred unchanged to radially populated three-dimensional families, where the number of modes in each shell also grows. The operator domain is a condition on the complete weighted sum, not a visual smoothness judgment.


\clearpage
\section{A graph approximation controls the field and its image}
\atlasfigure{09_graph_truncation}{Truncation errors for the normalized positive-axis sequence $a_n=n^{-2}/\sqrt{\zeta(4)}$ at $m=1$. Both curves use infinite tail formulas, evaluated with the Hurwitz zeta function. Dotted lines show asymptotic slopes rather than fitted bounds.}

The graph norm is
\[
\norm{a}_{H_m}^2=\norm{a}^2+\norm{H_ma}^2.
\]
Convergence in ordinary $\ell^2$ controls the field. Graph convergence also controls the result of applying the unbounded operator. For the displayed sequence,
\[
\norm{a-P_Ka}^2=\frac{\zeta(4,K+1)}{\zeta(4)},\qquad
\norm{a-P_Ka}_{H_1}^2=
\frac{2\zeta(4,K+1)+4\pi^2\zeta(2,K+1)}{\zeta(4)}.
\]
The ordinary error decays like $K^{-3/2}$ while the graph error decays like $K^{-1/2}$. Both vanish, but the unbounded multiplier makes high modes more expensive in the second measurement.

The maximal-domain argument proves the structural statement behind this example: finite truncations preserve the domain, commute with the multiplier, and approximate a domain vector together with its image. This gives a finite Fourier graph core. No observed slope is needed for that existence statement.


\clearpage
\section{Fourier transport carries the domain as well as the formula}
\atlasfigure{10_fourier_transport}{The inverse Fourier equivalence carries coefficient fields to physical torus fields. The square commutes by the definition of the transported operator. The upper horizontal map acts on the exact domains, and the lower map acts on the ambient Hilbert spaces.}

Let $\mathcal F$ send a physical field to its Fourier coefficients. The physical Hamiltonian is defined by
\[
\mathsf H_m=\mathcal F^{-1}H_m\mathcal F,
\qquad
D(\mathsf H_m)=\mathcal F^{-1}D(H_m).
\]
The domain clause is essential. The lower horizontal equivalence is defined everywhere on the Hilbert space, but the vertical arrows need not be. They are unbounded operators with specified dense domains.

Because Fourier transport is unitary, it preserves inner products and transfers self-adjointness and closedness. It also gives the physical form of the weighted domain condition:
\[
f\in D(\mathsf H_m)
\iff \sum_{n\in\Z^3}(1+|2\pi n|^2)\norm{\widehat f(n)}^2<\infty.
\]
This is a genuine physical-space operator on torus spinor fields. The square does not by itself identify a previously constructed time-evolution group with the generator of this operator. That relationship is proved in the generator chapter by a full-domain argument, rather than by the diagram alone.


\clearpage
\section{Changing mass changes the action, not the domain}
\atlasfigure{11_mass_domain}{Left: the graph weight relative to the first-order Fourier weight for three masses. Right: the exact graph norm square of one fixed normalized domain vector. The physical mass interpretation often takes $m\geq0$, while the mathematics permits every real $m$.}

For a square-summable coefficient field, the energy identity is
\[
\norm{H_ma}^2=m^2\norm{a}^2+
\sum_n|2\pi n|^2\norm{a_n}^2.
\]
The mass contribution is finite whenever $a\in\ell^2$. Finiteness of the unbounded part therefore depends only on the momentum-weighted sum. Consequently $D(H_{m_1})=D(H_{m_2})$ for any two real masses.

Equivalent graph norms need not have equal numerical values. For the normalized axis family $a_n=n^{-2}/\sqrt{\zeta(4)}$, the momentum contribution is exactly
\[
\frac{4\pi^2\zeta(2)}{\zeta(4)}=60.
\]
The field norm is one, hence $\norm{a}_{H_m}^2=61+m^2$. A larger mass raises this graph norm while leaving the vector inside the same domain. The graph-weight ratios in the left panel approach one at high momentum because the derivative term dominates.


\clearpage
\section{One construction, distinct evidence layers}
\atlasfigure{12_frontier}{The checked operator, generator, centered limits and third-order moment estimates, the separate spectral-calculus API linkage, and the comparisons required by an interacting or sampled-grid model. The layers state verification scope rather than physical completeness.}

Finite matrices and finite sums supply exact local identities. Completion removes the Fourier cutoff. Fourier transport identifies physical fields, and the adjoint-domain calculation determines the unbounded Hamiltonian. The checked generator theorem then ties the completed group to that operator on the full domain and gives the strong Schrödinger equation. The abstract spectral-functional-calculus identification remains a separate written argument.

For the fixed centered Fourier family, the checked arguments establish graph-domain and strong-resolvent convergence as spacing approaches zero through either sign. A finite third-order spatial Fourier moment gives the quadratic compact-time approximation rate. The technical companion identifies the actual declarations, domain hypotheses, and accepted verification record.

A changing finite spatial lattice needs sampling, interpolation, aliasing control, and compatible geometry. Background coupling changes the operator-domain problem. Dynamical gauge and gravitational fields require additional evolution and constraint arguments. The diagram does not identify those comparisons with the free-field theorem.

\clearpage
\section{Reading the plates as one argument}

The first three plates fix the representation. Refinement resolves an operator into sectors, normalized Fourier modes identify its spatial channels, and the Clifford matrices fix how each channel acts on a four-component spinor. No convergence claim is needed at this stage. The central facts are exact algebraic identities and exact normalization choices.

The next three plates introduce evolution and approximation. The sine symbol changes the dispersion relation while retaining Hermiticity. Its small-spacing error is quadratic at fixed momentum. Exact modal unitarity then protects the total packet norm even when the local density moves. Stability and consistency serve different purposes: stability prevents growth of an error that is already present, while consistency makes the finite-core error small as the spacing decreases.

The two-tail plate explains how these finite facts reach a full Hilbert space. First approximate a given vector by finite data. Then refine the spacing on that fixed finite approximation. Unitarity controls the discarded part. The argument is powerful precisely because it does not require every square-summable vector to have finite derivative energy.

The remaining operator plates ask a different question: which fields can the continuum Hamiltonian itself act on as an $L^2$ operator? The weighted-domain criterion answers this exactly. Graph approximation strengthens ordinary approximation by controlling both a field and its image. Fourier transport carries the resulting closed, self-adjoint operator and its domain to physical torus fields. The mass identity shows that this domain reflects spatial regularity rather than a choice of mass.

\paragraph{What the companion can establish.}
The atlas makes the source architecture inspectable and supplies reproducible examples. The symbolic matrix checks are exact in their finite algebraic setting. The plotted calculations are bounded numerical evaluations with recorded conventions. Infinite-dimensional conclusions depend on the corresponding mathematical arguments and formal source statements.

\paragraph{What remains separate.}
The figures illustrate the mathematical argument without extending its verification scope. The generator and centered-approximation theorems have written proofs and corresponding formal declarations. The abstract spectral-calculus API linkage remains separate. None of these results supplies the coupled gravitational dynamics or the controlled production history required by the physical matter model.

\endgroup
